\section{Metodologia}\label{sec:methodology}

O protocolo separa três operações que costumam aparecer misturadas: a obtenção da resistência característica, que define a classe atribuída; a substituição da rigidez experimental pelo valor de classe; e o redimensionamento, que altera a geometria e pode compensar ou ampliar a diferença observada sob geometria fixa. Cada comparação mantém comuns todas as hipóteses que não constituem o objeto do experimento, de modo que a diferença medida seja atribuível à representação das propriedades dentro do modelo da Seção~\ref{sec:bridge_design}.

\subsection{Módulo de elasticidade comparado}\label{sec:modulo}

A fonte experimental publica dois módulos por espécie: o módulo na compressão paralela às fibras, $E_{c0}$, e o módulo na flexão estática, $E_{M0}$. A ABNT NBR 7190-3:2022 tabela apenas $E_{c0,med}$ por classe de resistência. A comparação adota $E_{c0}$ nos dois lados de cada par, com o valor medido da espécie no cenário experimental e o valor da classe no cenário normativo. Essa escolha mantém a mesma propriedade nos dois membros e dispensa um fator de conversão entre módulos, que seria uma hipótese adicional introduzida apenas no lado normativo.

Nas 40 espécies da base, a razão $E_{M0}/E_{c0}$ tem média de 0,97 e varia entre 0,73 e 1,13. A proximidade das médias sustenta o uso de um único módulo longitudinal no modelo de viga; a amplitude por espécie justifica não tratar os dois módulos como intercambiáveis caso a caso. O efeito de adotar $E_{M0}$ no lugar de $E_{c0}$ é reportado como análise de sensibilidade, e não como resultado principal.

\subsection{Cenários de propriedade}\label{sec:cenarios}

Para cada espécie $s$ e vão $L$, definem-se três cenários de propriedade sobre o mesmo modelo estrutural e o mesmo carregamento, conforme a Tabela~\ref{tab:cenarios}. O cenário \emph{esp} reúne as propriedades medidas da espécie e é o membro de referência dos dois experimentos. O cenário \emph{rig} difere de \emph{esp} apenas no módulo de elasticidade, substituído pelo da classe. O cenário \emph{cls} substitui o conjunto completo de propriedades pelos valores da classe.

\begin{table}[!ht]
\centering
\caption{Cenários de propriedade de cada par espécie--vão.}
\label{tab:cenarios}
\footnotesize
\begin{tabular}{llll}
\toprule
Propriedade & \emph{esp} & \emph{rig} & \emph{cls} \\
\midrule
Módulo de elasticidade, $E_{c0}$ & medido & da classe & da classe \\
Densidade aparente, $\rho_{ap}$ & medida & medida & da classe \\
Resistência à flexão, $f_{m,k}$ & medida & medida & da classe \\
Resistência ao cisalhamento, $f_{v0,k}$ & medida & medida & da classe \\
\bottomrule
\end{tabular}
\end{table}

Desses cenários decorrem dois experimentos. O experimento principal, \emph{esp} contra \emph{rig}, é controlado: como só o módulo muda, a diferença observada é atribuível à representação da rigidez. O experimento secundário, \emph{esp} contra \emph{cls}, reproduz a decisão de projeto efetivamente tomada quando se dimensiona pela classe, e nele o efeito da rigidez aparece combinado com os das demais propriedades. Reportar os dois separa o mecanismo da sua consequência prática, e a comparação entre eles indica em que medida as demais propriedades compensam ou reforçam o efeito da rigidez.

A resistência característica à flexão adota $f_{m,k}=f_{c0,k}$ nos dois lados, conforme o item 6.3.4 da \textcite{NBR7190-1:2022}. A espécie entra com o $f_{c0,k}$ publicado por \textcite{DiasLahr2004}, obtido pelo estimador de estatística de ordem da norma, e não por conversão simplificada a partir da média. A resistência característica ao cisalhamento da espécie é estimada a partir da média publicada, única forma disponível na fonte, e está marcada como valor estimado na base de dados; essa entrada é comum aos cenários \emph{esp} e \emph{rig} e, portanto, não interfere no experimento principal.

\subsection{Comparação sob geometria fixa}\label{sec:geometria_fixa}

Antes do redimensionamento, a mesma estrutura é avaliada com as duas rigidezes, o que isola a resposta do modelo da resposta do otimizador. O erro de previsão do deslocamento é
\begin{equation}
\varepsilon_{\delta,s,L}=\left(\frac{\delta_C}{\delta_R}-1\right)100\%.
\label{eq:erro_delta}
\end{equation}
No regime elástico linear, com ações e geometria fixas e o deslocamento inversamente proporcional a um único módulo,
\begin{equation}
\frac{\delta_C}{\delta_R}=\frac{E_{c0,s}}{E_{c0,c_s}},
\label{eq:razao_delta}
\end{equation}
em que $E_{c0,c_s}$ é o módulo da classe $c_s$ atribuída à espécie $s$. A Equação~\eqref{eq:razao_delta} serve como verificação analítica fechada da implementação sob essas condições, e não deve ser estendida a modelos com parcelas de deformação independentes desse módulo.

\subsection{Redimensionamento}\label{sec:optimisation}

Cada cenário é redimensionado por otimização, sob critérios de aceitação idênticos. O vetor de projeto é $\mathbf{x}=(d,b_w,h,esp_{long},esp_{tab})^{T}$, com os intervalos de busca da Tabela~\ref{tab:intervalos}, comuns a todos os casos.

\begin{table}[!ht]
\centering
\caption{Intervalos de busca das variáveis de projeto.}
\label{tab:intervalos}
\footnotesize
\begin{tabular}{lcc}
\toprule
Variável & Mínimo (cm) & Máximo (cm) \\
\midrule
Diâmetro da longarina, $d$ & 20 & 100 \\
Largura da prancha do tabuleiro, $b_w$ & 20 & 50 \\
Altura da prancha do tabuleiro, $h$ & 5 & 15 \\
Espaçamento entre longarinas, $esp_{long}$ & 30 & 200 \\
Espaçamento entre peças do tabuleiro, $esp_{tab}$ & 2 & 5 \\
\bottomrule
\end{tabular}
\end{table}

O problema é biobjetivo, minimizando o volume de madeira da Equação~\eqref{eq:volume} e a utilização do limite de serviço, sujeito às seis restrições da Seção~\ref{sec:restricoes}:
\begin{equation}
\min_{\mathbf{x}\in\mathcal{D}}\left(V(\mathbf{x}),\,U_{fin}(\mathbf{x})\right),
\qquad \widehat g_j(\mathbf{x})\leq0,\quad j=1,\ldots,6,
\label{eq:problema}
\end{equation}
em que $\mathcal{D}$ reúne os intervalos da Tabela~\ref{tab:intervalos} e a admissibilidade geométrica da Equação~\eqref{eq:acomodacao}.

A avaliação incorpora tolerância dimensional do diâmetro, que é a dimensão sujeita a maior desvio em peças roliças por ser resultado da seleção da tora e não de um processo de corte controlado. Os objetivos são avaliados na geometria nominal, porque volume e flecha reportados são os do projeto tal como especificado. As restrições são avaliadas no pior caso de uma grade determinística de $N_c=5$ níveis aplicada ao diâmetro,
\begin{equation}
\widetilde d^{\,(c)}=d\left(1+\rho\,\xi^{(c)}\right),\quad
\xi^{(c)}\in\left\{-1,-\tfrac{1}{2},0,+\tfrac{1}{2},+1\right\},\qquad
\widehat g_j(\mathbf{x})=\max_{c}\,g_j\!\left(\widetilde{\mathbf{x}}^{(c)}\right),
\label{eq:agregacao_robusta}
\end{equation}
com $\rho=5\%$ em toda a campanha. A solução só é aceita se sobreviver a todos os desvios do intervalo, e não em média sobre eles. A grade é determinística e idêntica em todos os cenários, o que preserva o pareamento: dois cenários de um mesmo par enfrentam exatamente o mesmo conjunto de perturbações. Como a acomodação geométrica da Equação~\eqref{eq:acomodacao} é recalculada a cada perturbação, o número inteiro de peças pode mudar dentro da grade; a formulação representa, portanto, sensibilidade do processo de dimensionamento, e não apenas tolerância de fabricação de uma ponte com número de peças fixo.

O valor $\rho=5\%$ é mantido fixo. A campanha não varre níveis de robustez, porque o objeto do estudo é a representação das propriedades do material, e não a tolerância geométrica; fixar $\rho$ garante que essa hipótese seja idêntica nos dois membros de cada par e não participe da diferença medida.

A busca usa o NSGA-II \parencite{Deb2002} implementado em \texttt{pymoo} \parencite{blank2020}, com a configuração da Tabela~\ref{tab:otimizacao}. De cada fronteira extrai-se a solução viável de menor volume, que é o extremo de interesse para o anteprojeto e o ponto em que a comparação de consumo é mais exigente.

\begin{table}[!ht]
\centering
\caption{Configuração do NSGA-II, comum a todos os casos.}
\label{tab:otimizacao}
\footnotesize
\begin{tabular}{lc}
\toprule
Parâmetro & Valor \\
\midrule
População / gerações & 50 / 300 \\
Inicialização & Aleatória uniforme \\
Cruzamento SBX & $p_c=0{,}90$, $\eta_c=15$ \\
Mutação polinomial & $\eta_m=20$ \\
Eliminação de duplicatas & Sim \\
Níveis da grade de robustez, $N_c$ & 5 \\
Amplitude relativa, $\rho$ & 5\% \\
Sementes pareadas por cenário & 5 \\
\bottomrule
\end{tabular}
\end{table}

\subsection{Repetição e separação entre efeito e ruído}\label{sec:sementes}

Cada cenário é executado com sementes pareadas: os três cenários de um mesmo par espécie--vão compartilham a semente, e o conjunto é repetido para sementes distintas. Essa repetição é necessária porque uma execução única por condição não distingue um efeito de material da dispersão do próprio otimizador entre execuções. A diferença de volume é calculada par a par, dentro de cada semente, e só então agregada entre sementes, o que preserva o pareamento e permite reportar a dispersão da diferença junto com o seu valor central. Uma diferença cuja magnitude seja da ordem da dispersão entre sementes não será interpretada como efeito de material.

\subsection{Diferença de volume e reavaliação do projeto}\label{sec:protocolo}

Para cada par, sejam $\mathbf{x}_R$ e $\mathbf{x}_C$ os projetos obtidos nos cenários experimental e normativo. A diferença relativa de volume é
\begin{equation}
\Delta V_{s,L}=\frac{V(\mathbf{x}_C)-V(\mathbf{x}_R)}{V(\mathbf{x}_R)}100\%.
\label{eq:deltaV}
\end{equation}
Um valor negativo indica que o projeto pela classe consome menos madeira, o que não constitui, por si só, evidência de violação de limite: pode decorrer de a classe atribuir à espécie uma rigidez superior à medida, caso em que o projeto é menos conservador, ou de a verificação governante não ser sensível à rigidez.

Para separar consumo de desempenho, a geometria do projeto normativo é mantida e recalculada com as propriedades medidas da espécie, sem reotimizar:
\begin{equation}
U_{fin,s,L}^{C\rightarrow R}=\frac{\delta_{fin}(\mathbf{x}_C;E_{c0,s})}{L/350},
\qquad
U_{inst,s,L}^{C\rightarrow R}=\frac{\delta_{inst}(\mathbf{x}_C;E_{c0,s})}{L/500}.
\label{eq:reavaliacao}
\end{equation}
As duas utilizações são reportadas, pois um projeto pode atender a uma combinação e exceder a outra. Valores superiores a um indicam excedência do respectivo limite de serviço nessa reavaliação, sob as hipóteses de ações, fluência e propriedades adotadas. Esse resultado é uma frequência de excedência no conjunto analisado, e não uma probabilidade de falha nem demonstração de insegurança no estado limite último. As demais verificações e a admissibilidade geométrica são igualmente reavaliadas, distinguindo avaliação nominal e avaliação sob a grade de perturbação.

\subsection{Matriz da campanha}\label{sec:campanha}

A campanha cobre as 40 espécies da base, quatro vãos e os três cenários, repetidos com cinco sementes pareadas, totalizando 2400 execuções independentes do NSGA-II. Os vãos adotados são $L=3$, 5, 8 e 10~m, com largura de pista $B=4{,}5$~m, seção transversal usual de ponte vicinal de pista simples. Toda a faixa permanece no primeiro trecho do coeficiente de impacto e abaixo do limite de $10$~m, de modo que nenhuma mudança de expressão de carregamento atravessa a faixa comparada e possa ser confundida com efeito de material.

A verificação governante é identificada em cada solução pelas utilizações e folgas normalizadas, registrando restrições simultaneamente ativas. Não se pressupõe que a flecha governe nem que ocorra migração entre estados limites ao longo da faixa de vãos. Se o deslocamento governar em toda a faixa, a análise permanece centrada na magnitude das diferenças; se outra restrição limitar o volume, examina-se quanto isso reduz a sensibilidade do dimensionamento à rigidez, que é a condição sob a qual a representação por classe deixa de ser determinante.

\subsection{Verificação da implementação}\label{sec:verificacao}

A verificação independente compreende: a conferência da resposta sob geometria fixa contra a Equação~\eqref{eq:razao_delta}; a conferência das envoltórias das Equações~\eqref{eq:mqkenv} e~\eqref{eq:qqkenv} nos limites dos trechos de carregamento, em especial nos vãos curtos em que o diâmetro não é desprezível frente ao vão; a demonstração manual de um caso completo, reproduzida no Apêndice~\ref{ap:relatorio}; a repetição entre sementes; e a reavaliação das soluções selecionadas em sua geometria nominal. Os dados derivados e os scripts que geram a base experimental, a matriz de casos e as figuras são disponibilizados junto com o artigo.
